%% ============================================================================
%% Fichier  : annexes/Ann_A_glossaire.tex
%% Titre    : Glossaire
%% Version  : 1.0 -- novembre 2026 (INF3226, d'après ACSI)
%% ============================================================================
%% Définitions reprises des encadrés « Définition » des chapitres (termes
%% marqués \termecle), abrégées ; le chapitre de référence est indiqué.
%% ============================================================================

\chapter{Glossaire}
\label{ann:glossaire}

Ce glossaire reprend, sous une forme abrégée, les définitions données dans
les chapitres. Le numéro renvoie au chapitre où le terme est défini et
illustré~; c'est là qu'il faut lire la définition complète.

\begin{description}
    \item[\NEK{} (\textit{Not Enough Knowledge})] Réponse « connaissance
          insuffisante », réservée à un contrôle applicable dont on ne peut
          pas établir l'état~; ni \NON{}, ni sans objet
          (Ch.~\ref{chap:preuve}, \ref{chap:incertitude}).
    \item[ALE, SLE, ARO] Perte annuelle attendue, produit de la perte par
          occurrence et du taux annuel d'occurrence (Ch.~\ref{chap:risque}).
    \item[Appétence au risque] Niveau de risque que l'organisation accepte
          de prendre pour atteindre ses objectifs~; fixée par l'organe de
          gouvernance (Ch.~\ref{chap:risque}).
    \item[Assurance limitée] Niveau de confiance plus faible, exprimé sous
          une forme négative, après des travaux moins étendus
          (Ch.~\ref{chap:audit-auditologie}).
    \item[Assurance raisonnable] Niveau élevé, mais non absolu, de confiance,
          exprimé sous une forme positive, après des travaux étendus
          (Ch.~\ref{chap:audit-auditologie}).
    \item[Audit de conformité] Selon l'ISSAI~400, audit qui évalue si les
          activités, transactions et informations d'une entité publique sont
          conformes aux autorités qui la régissent~; relation à trois
          parties (Ch.~\ref{chap:normes}).
    \item[Auditologie] Science appliquée de la réduction de l'incertitude
          sur les missions auditologiques (A0~; Ch.~\ref{chap:auditologie}).
    \item[Benford (loi de)] Distribution attendue des premiers chiffres
          significatifs de nombreuses séries comptables~; son écart signale
          une anomalie à expliquer, non une fraude (Ch.~\ref{chap:fraude}).
    \item[Cartographie des risques] Représentation des risques selon leur
          vraisemblance et leur impact, qui oriente l'effort d'audit
          (Ch.~\ref{chap:risque}).
    \item[Certification] Attestation, par un organisme tiers accrédité,
          qu'un système de management est conforme à une norme, sur un
          périmètre et pour une durée donnés (Ch.~\ref{chap:audit-auditologie}).
    \item[Chaîne de conservation] Traçabilité ininterrompue d'une pièce~:
          qui l'a détenue, quand, et pour quoi faire (Ch.~\ref{chap:audit-preuve}).
    \item[Charte d'audit] Document approuvé par l'organe de gouvernance, qui
          fixe la mission, les pouvoirs, le rattachement et l'indépendance
          de l'audit (Ch.~\ref{chap:mission}).
    \item[Conflit d'intérêts] Situation où un intérêt personnel ou une
          relation peut influencer, ou paraître influencer, le jugement de
          l'auditeur (Ch.~\ref{chap:deontologie}).
    \item[Conseil] Accompagnement qui aide à concevoir ou améliorer un
          dispositif~; le conseiller ne donne pas d'opinion indépendante sur
          ce qu'il a conçu (Ch.~\ref{chap:audit-auditologie}).
    \item[Constatation] Fait vérifiable, établi par l'auditeur et appuyé
          sur une preuve conservée (Prologue).
    \item[Contenant et contenu] Pour un document de gouvernance~: la manière
          dont il a été élaboré et validé, puis ce qu'il dit
          (Ch.~\ref{chap:audit-gouvernance}).
    \item[Contradictoire] Principe selon lequel l'audité peut répondre aux
          constats avant le rapport~; il corrige, rend opposable et réduit
          l'incertitude (Ch.~\ref{chap:mission}).
    \item[Contrôle continu, audit continu] Surveillance automatisée par la
          direction~; évaluation automatisée et fréquente par l'audit
          (Ch.~\ref{chap:audit-continu}).
    \item[Contrôle interne] Processus mis en œuvre par l'organe de
          gouvernance, la direction et le personnel, destiné à fournir une
          assurance raisonnable sur les objectifs d'opérations,
          d'information et de conformité (COSO~; Ch.~\ref{chap:controle-interne}).
    \item[Contrôle préventif, détectif, correctif] Contrôle qui empêche,
          repère ou corrige une anomalie (Ch.~\ref{chap:controle-interne}).
    \item[Contrôles d'application] Contrôles intégrés à une application,
          sur les entrées, les traitements et les sorties
          (Ch.~\ref{chap:controle-interne}).
    \item[Contrôles généraux informatiques] Contrôles de l'environnement
          des applications~: accès, changements, exploitation
          (Ch.~\ref{chap:controle-interne}).
    \item[CVSS] Système commun de notation de la gravité des
          vulnérabilités, de 0 à 10 (Ch.~\ref{chap:intrusion}).
    \item[Déclaration d'applicabilité] Document d'un SMSI qui liste les
          mesures de l'annexe~A de l'ISO/IEC~27001, retenues ou exclues,
          avec leur justification (Ch.~\ref{chap:audit-smsi}).
    \item[Décomposition] Organisation de l'espace d'audit en niveaux
          emboîtés, du domaine à la trace, sous la borne de sept
          (Ch.~\ref{chap:decomposition}).
    \item[Déficience majeure] Selon le COSO, déficience qui réduit fortement
          la probabilité d'atteindre les objectifs~; elle empêche de
          conclure à l'efficacité (Ch.~\ref{chap:referentiels}).
    \item[Diamant de la fraude] Extension du triangle de la fraude par la
          capacité de l'auteur (Wolfe et Hermanson~; Ch.~\ref{chap:fraude}).
    \item[Divergence de Kullback-Leibler] Mesure de l'écart entre les
          distributions de réponses de deux cycles d'audit
          (Ch.~\ref{chap:conclure}).
    \item[EBIOS Risk Manager] Méthode française d'appréciation des risques
          numériques, en cinq ateliers (Ch.~\ref{chap:gestion-risques}).
    \item[Entropie] Mesure de la dispersion des réponses~:
          $H = -\sum p \log_2 p$, au plus $\log_2 3$ pour trois modalités
          (Ch.~\ref{chap:conclure}).
    \item[Erreur] Inexactitude non intentionnelle (Ch.~\ref{chap:fraude}).
    \item[Espace d'audit] Ensemble des objets sur lesquels porte une
          mission, tel que le délimite la lettre de mission
          (Ch.~\ref{chap:decomposition}).
    \item[Expert judiciaire] Personne désignée par le juge pour répondre à
          des questions techniques~; distincte de l'auditeur et de
          l'enquêteur (Ch.~\ref{chap:audit-preuve}).
    \item[FAIR] Modèle quantitatif du risque qui décompose la perte en
          fréquence et en ampleur, souvent par simulation
          (Ch.~\ref{chap:gestion-risques}).
    \item[Fiche d'assertion] Fiche attachée à un contrôle~: objectif,
          critère, question, risque, conséquence (Ch.~\ref{chap:decomposition}).
    \item[Fraude] Acte intentionnel qui recourt à la tromperie pour obtenir
          un avantage indu ou causer un préjudice (Ch.~\ref{chap:fraude}).
    \item[Grille adaptative] Grille dont les questions suivantes dépendent
          des réponses obtenues~; l'effort va vers l'incertitude
          (Ch.~\ref{chap:audit-continu}).
    \item[Impression] Jugement formé sans démarche ni preuve~; elle oriente,
          elle n'entre jamais dans un rapport (Prologue).
    \item[Incertitude de type~A, de type~B] Selon le GUM, incertitude
          évaluée par des méthodes statistiques, ou par d'autres moyens
          (jugement, convention) (Ch.~\ref{chap:incertitude}).
    \item[Inspection] Vérification ponctuelle de conformité, souvent par une
          autorité de tutelle (Ch.~\ref{chap:audit-auditologie}).
    \item[Instrument trichotomique] Grille dont chaque question admet
          \OUI{}, \NON{} ou \NEK{} (A1~; Ch.~\ref{chap:preuve}).
    \item[Intervalle de Wilson] Intervalle de confiance d'une proportion
          observée, adapté aux petits effectifs (Ch.~\ref{chap:incertitude}).
    \item[Mission auditologique] Quintuplet
          $\langle \mathcal{S}, \mathcal{A}, \mathcal{I}, \mathcal{G},
          \mathcal{T} \rangle$~: système, architecture cognitive, canal
          d'information, structure stratégique, trajectoire
          (A0~; Ch.~\ref{chap:auditologie}).
    \item[Modèle de responsabilité partagée] Répartition des contrôles
          entre le fournisseur de nuage et son client, selon le modèle de
          service (Ch.~\ref{chap:audit-tiers}).
    \item[Modèle des trois lignes] Répartition des rôles entre organe de
          gouvernance, management et audit interne (IIA~;
          Ch.~\ref{chap:audit-auditologie}).
    \item[Norme professionnelle d'audit] Exigences qui encadrent la conduite
          de l'audit, la compétence et l'éthique de l'auditeur, et la forme
          des conclusions (Ch.~\ref{chap:normes}).
    \item[Opinion d'audit] Conclusion motivée tirée d'un ensemble de
          constatations, au regard d'un critère connu à l'avance
          (Prologue).
    \item[Papiers de travail] Documents qui retracent les travaux, les
          preuves et le raisonnement de l'auditeur (Ch.~\ref{chap:preuve}).
    \item[Planification] Phase qui prépare le terrain à partir de la lettre
          de mission, jusqu'au plan d'audit (Ch.~\ref{chap:planifier}).
    \item[Preuve d'audit] Information, suffisante et appropriée, sur
          laquelle l'auditeur fonde sa cotation (Ch.~\ref{chap:preuve}).
    \item[Prise de connaissance] Compréhension, avant le terrain, de
          l'organisation, de son SI et de ses risques (Ch.~\ref{chap:planifier}).
    \item[Rapport ISAE~3402] Rapport d'assurance sur les contrôles d'une
          société de services~; le type~1 porte sur la conception à une
          date, le type~2 aussi sur l'efficacité sur une période
          (Ch.~\ref{chap:normes}).
    \item[Registre des traitements] Inventaire des traitements de données
          personnelles d'un responsable de traitement~: finalités,
          données, destinataires, durées, sécurité
          (Ch.~\ref{chap:audit-conformite}).
    \item[Règle des quatre yeux] Toute opération sensible est réalisée par
          une personne et validée par une autre (Ch.~\ref{chap:controle-interne}).
    \item[Risque inhérent, de contrôle, de détection] Risque d'anomalie avant
          contrôle~; risque que le contrôle interne ne l'arrête pas~; risque
          que l'auditeur ne la détecte pas. Leur produit est le risque
          d'audit (Ch.~\ref{chap:risque}).
    \item[Risque résiduel] Risque qui subsiste après l'effet des contrôles
          efficaces (Ch.~\ref{chap:risque}).
    \item[Séparation des tâches] Répartition, entre des personnes
          différentes, des fonctions dont le cumul permettrait de commettre
          et de dissimuler une erreur ou une fraude~: autoriser,
          enregistrer, garder l'actif, contrôler
          (Ch.~\ref{chap:controle-interne}).
    \item[Sondage en unités monétaires] Échantillonnage où chaque franc a
          la même chance d'être tiré~; il projette une erreur en montant
          (Ch.~\ref{chap:preuve}).
    \item[SUI] Indice d'usage stratégique de \NEK{}~: coefficient de l'enjeu
          des questions dans un modèle propre à chaque répondant~; indicateur
          de population, jamais de faute individuelle (A1~;
          Ch.~\ref{chap:fraude}, \ref{chap:incertitude}).
    \item[Taux d'incertitude ($U$)] Part des réponses \NEK{}
          (Ch.~\ref{chap:conclure}).
    \item[Taux de conformité ($\mu$)] Part des réponses \OUI{}
          (Ch.~\ref{chap:conclure}).
    \item[Test d'intrusion] Évaluation autorisée de la sécurité d'un
          système par la simulation d'attaques, dans un périmètre et une
          fenêtre convenus par écrit (Ch.~\ref{chap:intrusion}).
    \item[Travaux préparatoires] Phase qui conduit à la lettre de mission et
          à la constitution de l'équipe (Ch.~\ref{chap:planifier}).
    \item[Triangle de la fraude] Pression, opportunité, rationalisation~:
          les trois conditions de la fraude selon Cressey
          (Ch.~\ref{chap:fraude}).
    \item[Vraisemblance, impact, menace, vulnérabilité] Chance qu'un événement
          survienne~; ampleur de ses conséquences~; cause potentielle d'un
          incident~; faiblesse qu'elle peut exploiter (Ch.~\ref{chap:risque}).
\end{description}
